% ============================================================================
%  ESTADÍSTICA DESCRIPTIVA — Sesión 10
%  Semana 10: 12 al 17 de octubre de 2026   (lunes 12: festivo)
%  2.6 Interpretación de resultados: patrones y tendencias
%  2.7 Redacción de informes estadísticos
%  Duración prevista: 3 horas
% ============================================================================
\documentclass[aspectratio=169,10pt]{beamer}

\usepackage[utf8]{inputenc}
\IfFileExists{spanish.ldf}{\usepackage[spanish,es-noshorthands]{babel}}{}
\usepackage{amsmath,amssymb}
\usepackage{booktabs}
\usepackage{tikz}
\usepackage{listings}
\usepackage{pgfplots}
\pgfplotsset{compat=1.16}

\usetheme{Uniminuto}

\uniasignatura{Estadística descriptiva}
\uniprofesor{Andrés Felipe Puerta Velez}
\unicorreo{andres.puerta-v@uniminuto.edu.co}
\uninrc{CBASBL031 - CBASBL141}
\unisemana{Sesión 10 — Semana 10 · 12 al 17 de octubre de 2026}

\newcommand{\listatiny}{\setbeamerfont{itemize/enumerate body}{size=\tiny}\tiny}

\pgfplotsset{
  uniplot/.style={
    width=8.0cm, height=3.5cm, axis lines=left,
    tick label style={font=\tiny,color=uniGris},
    label style={font=\tiny,color=uniGris},
    title style={font=\scriptsize\bfseries,color=uniAzul},
    every axis plot/.append style={fill=uniAzul!70,draw=uniAzul},
    ymajorgrids, grid style={uniGris!25,very thin},
  }
}

\newcommand{\fx}[1]{{\ttfamily\color{uniAzul}#1}}

% ------------------------------------------------------------- CÓDIGO R ----
\lstdefinestyle{Rcode}{
  language=R,
  basicstyle=\ttfamily\scriptsize,
  keywordstyle=\color{uniAzul}\bfseries,
  commentstyle=\color{uniGris}\itshape,
  stringstyle=\color{uniTitulo},
  showstringspaces=false, numbers=none, breaklines=true,
  columns=fullflexible,
  backgroundcolor=\color{uniCremaClaro},
  frame=leftline, framerule=1.2pt, rulecolor=\color{uniNaranja},
  xleftmargin=0.35cm, framexleftmargin=0.35cm,
  framexrightmargin=0.15cm, framextopmargin=2pt, framexbottommargin=2pt,
  aboveskip=0.15cm, belowskip=0.15cm
}
\lstdefinestyle{Rmini}{
  language=R,
  basicstyle=\ttfamily\tiny,
  keywordstyle=\color{uniAzul}\bfseries,
  commentstyle=\color{uniGris}\itshape,
  stringstyle=\color{uniTitulo},
  showstringspaces=false, numbers=none, breaklines=true,
  columns=fullflexible,
  backgroundcolor=\color{uniCremaClaro},
  frame=leftline, framerule=1.2pt, rulecolor=\color{uniNaranja},
  xleftmargin=0.35cm, framexleftmargin=0.35cm,
  framexrightmargin=0.15cm, framextopmargin=2pt, framexbottommargin=2pt,
  aboveskip=0.12cm, belowskip=0.12cm
}
\lstdefinestyle{Rout}{
  basicstyle=\ttfamily\tiny\color{black},
  showstringspaces=false, numbers=none, breaklines=true,
  columns=fullflexible,
  backgroundcolor=\color{uniGris!10},
  frame=leftline, framerule=1.2pt, rulecolor=\color{uniGris},
  xleftmargin=0.35cm, framexleftmargin=0.35cm,
  framexrightmargin=0.15cm, framextopmargin=2pt, framexbottommargin=2pt,
  aboveskip=0.12cm, belowskip=0.12cm
}

\begin{document}

\portada

\begin{frame}{Orden del día}
  \begin{enumerate}
    \item Llamado a lista y revisión de los ejercicios propuestos
    \item \azul{Tema 2.6: de la cifra a la afirmación}
    \item Patrones: concentración, dispersión, forma y atípicos
    \item Tendencias: comparar grupos y comparar en el tiempo
    \item Errores de interpretación que cuestan dinero
    \item \azul{Tema 2.7: redacción del informe estadístico} y taller final
  \end{enumerate}

  \vspace{0.25cm}
  {\footnotesize\color{uniGris} Sesión de 3 horas. Hoy no aprendemos fórmulas
  nuevas: aprendemos a \azul{usar} las que ya tenemos.}
\end{frame}

\begin{frame}{Todo lo que sabemos del Almacén Bello}
  \scriptsize
  \begin{center}
  \begin{tabular}{@{}lrlr@{}}
    \toprule
    \textbf{Medida} & \textbf{Valor} & \textbf{Medida} & \textbf{Valor} \\
    \midrule
    $n$ & 40 & Rango $R$ & 240 \\
    Media $\bar{x}$ & 227{,}00 & $RIC$ & 81{,}27 \\
    Mediana $Me$ & 224{,}73 & Desviación $s$ & 56{,}34 \\
    Moda $Mo$ & 218{,}67 & $CV$ & 24{,}8\% \\
    $Q_1$ & 186{,}29 & Asimetría $As_2$ & $+0{,}121$ \\
    $Q_3$ & 267{,}56 & Curtosis $g_2$ & $-0{,}741$ \\
    \bottomrule
  \end{tabular}
  \end{center}
  {\tiny\color{uniGris} Gasto mensual, en miles de pesos, de 40 clientes.
  Valores calculados con la tabla agrupada.}

  \vspace{0.1cm}
  \begin{alertblock}{La pregunta de hoy}
    \scriptsize
    Tenemos doce números. \destacar{¿Qué le decimos al gerente?} Una lista de
    cifras no es un análisis: el trabajo profesional consiste en convertirlas
    en tres o cuatro afirmaciones que sirvan para decidir.
  \end{alertblock}
\end{frame}

\seccion{2.6 Interpretación de resultados}

\begin{frame}{Anatomía de una buena interpretación}
  \footnotesize
  \begin{block}{Los cuatro elementos}
    \scriptsize
    \begin{enumerate}\setlength{\itemsep}{-0.03cm}
      \item \azul{Qué se midió} y sobre quién: ``el gasto mensual de 40
            clientes''.
      \item \azul{La cifra con su unidad}: ``227 mil pesos''.
      \item \azul{Qué significa} en lenguaje del negocio: ``es el gasto
            promedio''.
      \item \azul{Con qué se compara} o qué implica: ``pero la mitad gasta
            menos de 225 mil''.
    \end{enumerate}
  \end{block}

  \vspace{0.1cm}
  \begin{columns}[T,onlytextwidth]
    \begin{column}{0.48\textwidth}
      \begin{alertblock}{Mal}
        \tiny
        ``La media es 227 y la desviación 56,34. El CV es 24,8\% y la
        asimetría 0,121.''
      \end{alertblock}
    \end{column}
    \begin{column}{0.48\textwidth}
      \begin{exampleblock}{Bien}
        \tiny
        ``Los 40 clientes de la muestra gastan en promedio 227 mil pesos al
        mes, con una variación típica de 56 mil. La mitad gasta menos de 225
        mil, y el cuarto más gastador supera los 268 mil.''
      \end{exampleblock}
    \end{column}
  \end{columns}

  \vspace{0.05cm}
  {\tiny\color{uniGris} La segunda versión no tiene ni una fórmula, y contiene
  exactamente la misma información.}
\end{frame}

\begin{frame}{Frases plantilla para cada medida}
  \scriptsize
  \begin{center}
  \begin{tabular}{@{}p{2.4cm}p{9.4cm}@{}}
    \toprule
    \textbf{Medida} & \textbf{Cómo se enuncia} \\
    \midrule
    Media & ``En promedio, cada [unidad] registra [valor] [unidad de medida].'' \\
    Mediana & ``La mitad de los [casos] está por debajo de [valor].'' \\
    Moda & ``El valor más frecuente es [valor], con [$f_i$] casos.'' \\
    $Q_1$, $Q_3$ & ``El 25\% más bajo no supera [valor]; el 25\% más alto
      supera [valor].'' \\
    Desviación & ``Los datos se apartan del promedio, en promedio,
      [valor] [unidad].'' \\
    $CV$ & ``El grupo es [homogéneo / heterogéneo]: la variación equivale al
      [valor]\% de la media.'' \\
    Asimetría & ``La distribución tiene una cola hacia [la derecha / la
      izquierda]: unos pocos casos [muy altos / muy bajos].'' \\
    Atípicos & ``[n] casos se salen del patrón general y conviene revisarlos
      aparte.'' \\
    \bottomrule
  \end{tabular}
  \end{center}
\end{frame}

\begin{frame}{Patrones: los cinco que siempre se buscan}
  \scriptsize
  \begin{center}
  \begin{tabular}{@{}p{2.5cm}p{4.4cm}p{5.0cm}@{}}
    \toprule
    \textbf{Patrón} & \textbf{Cómo se detecta} & \textbf{Qué suele significar} \\
    \midrule
    Concentración & $CV$ bajo, $RIC$ estrecho, histograma con pico &
      Grupo homogéneo; se puede tratar con una sola estrategia \\
    \addlinespace
    Dispersión & $CV$ alto, bigotes largos &
      Hay subgrupos distintos escondidos en el promedio \\
    \addlinespace
    Asimetría & $\bar{x} \ne Me$, un bigote más largo &
      Existe una minoría que jala el promedio \\
    \addlinespace
    Atípicos & Puntos fuera de las vallas &
      Errores de registro o casos excepcionales que hay que estudiar \\
    \addlinespace
    Bimodalidad & Dos picos en el histograma &
      Dos poblaciones mezcladas: hay que separarlas \\
    \bottomrule
  \end{tabular}
  \end{center}

  \vspace{0.05cm}
  \begin{alertblock}{El patrón del Almacén Bello}
    \tiny
    $CV = 24{,}8\%$ (aceptable), $\bar{x}$ ligeramente mayor que $Me$
    (asimetría positiva leve), sin atípicos, un solo pico.
    \destacar{Clientela razonablemente homogénea con una minoría de gasto
    alto}: se puede diseñar una oferta principal y una secundaria.
  \end{alertblock}
\end{frame}

\begin{frame}{Tendencias: comparar en el tiempo}
  \footnotesize
  \begin{columns}[T,onlytextwidth]
    \begin{column}{0.50\textwidth}
      \centering
      \begin{tikzpicture}
        \begin{axis}[uniplot, width=6.4cm, height=3.2cm,
          title={Ventas mensuales (millones)},
          symbolic x coords={Ene,Feb,Mar,Abr,May,Jun},
          xtick=data, ymin=100, ymax=170,
          every axis plot/.append style={fill=none}]
          \addplot[color=uniAzul,thick,mark=*,mark size=1.4pt,
                   mark options={fill=uniAzul}]
            coordinates {(Ene,120) (Feb,128) (Mar,135) (Abr,132) (May,145)
                         (Jun,158)};
        \end{axis}
      \end{tikzpicture}
    \end{column}
    \begin{column}{0.47\textwidth}
      \centering\scriptsize
      \begin{tabular}{@{}lrr@{}}
        \toprule
        \textbf{Mes} & \textbf{Ventas} & \textbf{Var. \%} \\
        \midrule
        Enero   & 120 & --- \\
        Febrero & 128 & $+6{,}7$ \\
        Marzo   & 135 & $+5{,}5$ \\
        Abril   & 132 & $-2{,}2$ \\
        Mayo    & 145 & $+9{,}8$ \\
        Junio   & 158 & $+9{,}0$ \\
        \bottomrule
      \end{tabular}
    \end{column}
  \end{columns}

  \vspace{0.1cm}
  \begin{alertblock}{Cómo se calcula y cómo se lee}
    \scriptsize
    \[ \text{Variación \%} = \frac{V_{\text{actual}} - V_{\text{anterior}}}
       {V_{\text{anterior}}} \times 100 \qquad
       \text{Enero a junio: } \frac{158-120}{120}\times100 = \destacar{+31{,}7\%} \]
    \azul{Tendencia creciente} con una caída puntual en abril. El crecimiento
    acumulado del semestre es del 31,7\%.
  \end{alertblock}
\end{frame}

\begin{frame}{Tendencias: comparar grupos}
  \scriptsize
  \begin{center}
  \begin{tabular}{@{}lrrrr@{}}
    \toprule
    & \textbf{Bello} & \textbf{Itagüí} & \textbf{Diferencia} &
    \textbf{Lectura} \\
    \midrule
    $n$ & 40 & 38 & --- & Muestras comparables \\
    Media & 227{,}0 & 198{,}5 & $-28{,}5$ & Itagüí vende menos por cliente \\
    Mediana & 224{,}7 & 201{,}0 & $-23{,}7$ & La diferencia no viene de atípicos \\
    $s$ & 56{,}3 & 71{,}2 & $+14{,}9$ & Itagüí es más variable \\
    $CV$ & 24{,}8\% & 35{,}9\% & $+11{,}1$ & Clientela menos homogénea \\
    \bottomrule
  \end{tabular}
  \end{center}

  \vspace{0.1cm}
  \begin{alertblock}{Cómo se redacta esta comparación}
    \scriptsize
    ``El gasto promedio en Itagüí (198\,500) es 28\,500 pesos menor que en
    Bello, y la diferencia se mantiene en la mediana, de modo que
    \destacar{no se explica por unos pocos clientes extremos}. Además, la
    clientela de Itagüí es más heterogénea ($CV$ del 35,9\% frente al
    24,8\%), lo que sugiere revisar si allí conviven dos segmentos
    distintos.''
  \end{alertblock}
\end{frame}

\begin{frame}{Cuatro errores de interpretación que cuestan dinero}
  \scriptsize
  \begin{center}
  \begin{tabular}{@{}p{3.0cm}p{4.4cm}p{4.6cm}@{}}
    \toprule
    \textbf{Error} & \textbf{Ejemplo} & \textbf{Corrección} \\
    \midrule
    Confundir asociación con causa &
      ``Subieron las ventas cuando pusimos el aviso'' &
      Pudo ser la temporada, el precio o la competencia \\
    \addlinespace
    La falacia del promedio &
      ``El cliente promedio gasta 227 mil'' &
      Puede que \azul{ningún} cliente gaste eso; mire la distribución \\
    \addlinespace
    Muestra no representativa &
      Encuestar solo los sábados en la mañana &
      La conclusión describe a ese momento, no al negocio \\
    \addlinespace
    Extrapolar sin base &
      ``Creciendo 9\% mensual, en un año duplicamos'' &
      La tendencia pasada no garantiza la futura \\
    \bottomrule
  \end{tabular}
  \end{center}

  \vspace{0.05cm}
  \begin{alertblock}{La prueba que hay que aplicarse siempre}
    \tiny
    Antes de escribir una conclusión, pregúntese: \destacar{¿qué otra
    explicación tendría estos mismos datos?} Si se le ocurre una, la
    conclusión debe suavizarse o debe recogerse más información.
  \end{alertblock}
\end{frame}

\begin{frame}{Ejercicio 1 — en clase}
  \begin{alertblock}{En parejas (20 minutos)}
    \scriptsize
    Resultados de una encuesta sobre tiempo de entrega (días) en dos
    proveedores:

    \vspace{0.15cm}
    \centering
    \begin{tabular}{@{}lrr@{}}
      \toprule
      & \textbf{Proveedor A} & \textbf{Proveedor B} \\
      \midrule
      $n$ & 50 & 50 \\
      Media & 6{,}2 & 5{,}8 \\
      Mediana & 6{,}0 & 4{,}5 \\
      $s$ & 1{,}1 & 3{,}9 \\
      Máximo & 9 & 22 \\
      \bottomrule
    \end{tabular}

    \vspace{0.2cm}
    \raggedright
    \begin{enumerate}\setlength{\itemsep}{-0.07cm}
      \item Calcule el $CV$ de cada proveedor.
      \item ¿Cuál tiene mejor promedio? ¿Cuál es más confiable?
      \item ¿Qué sugiere la diferencia entre media y mediana en B?
      \item Redacte tres frases de interpretación para el jefe de compras.
      \item ¿Con cuál proveedor firmaría? Justifique con dos cifras.
    \end{enumerate}
  \end{alertblock}
\end{frame}

\begin{frame}{Ejercicio 1 — solución}
  \scriptsize
  \begin{columns}[T,onlytextwidth]
    \begin{column}{0.42\textwidth}
      \begin{exampleblock}{Puntos 1 a 3}
        \scriptsize
        \[ CV_A = \frac{1{,}1}{6{,}2}\times100 = \destacar{17{,}7\%} \]
        \[ CV_B = \frac{3{,}9}{5{,}8}\times100 = \destacar{67{,}2\%} \]
        B tiene mejor promedio (5,8 frente a 6,2), pero A es
        \azul{mucho más confiable}.

        \vspace{0.1cm}
        En B, $\bar{x} > Me$ (5,8 frente a 4,5) y el máximo es 22:
        \destacar{asimetría positiva fuerte}, con entregas ocasionales muy
        tardías.
      \end{exampleblock}
    \end{column}
    \begin{column}{0.55\textwidth}
      \begin{block}{Puntos 4 y 5: redacción}
        \scriptsize
        ``El proveedor B entrega en 4,5 días en la mitad de los casos, frente a
        6 días de A: es \azul{más rápido cuando cumple}.''

        \vspace{0.08cm}
        ``Sin embargo, la variación de B equivale al 67\% de su media y
        registra entregas de hasta 22 días.''

        \vspace{0.08cm}
        ``El proveedor A nunca superó los 9 días y su variación es del 18\%.''

        \vspace{0.1cm}
        \destacar{Recomendación:} A, si la producción no tolera retrasos; B,
        solo para pedidos sin urgencia.
      \end{block}
    \end{column}
  \end{columns}
\end{frame}

\seccion{2.7 Redacción del informe estadístico}

\begin{frame}{Estructura del informe}
  \scriptsize
  \begin{center}
  \begin{tabular}{@{}rp{3.1cm}p{7.0cm}@{}}
    \toprule
    & \textbf{Sección} & \textbf{Contenido y extensión sugerida} \\
    \midrule
    1 & Título y resumen & Qué se estudió y el hallazgo principal (5 líneas) \\
    2 & Objetivo & La decisión que se quiere apoyar (2 líneas) \\
    3 & Metodología & Población, muestra, $n$, periodo, instrumento (1 párrafo) \\
    4 & Variables & Tabla con tipo y escala de cada una \\
    5 & Resultados & Tablas y gráficos numerados, con su lectura al pie \\
    6 & Interpretación & Patrones, tendencias y comparaciones (2--3 párrafos) \\
    7 & Limitaciones & Qué no permite concluir este estudio \\
    8 & Recomendaciones & Acciones concretas, cada una anclada a una cifra \\
    9 & Referencias & Fuentes de datos y bibliografía \\
    \bottomrule
  \end{tabular}
  \end{center}

  \vspace{0.05cm}
  {\tiny\color{uniGris} El resumen se escribe \azul{de último} aunque vaya de
  primero: solo se puede resumir lo que ya está escrito.}
\end{frame}

\begin{frame}{Cómo se presentan tablas y gráficos}
  \scriptsize
  \begin{columns}[T,onlytextwidth]
    \begin{column}{0.48\textwidth}
      \begin{block}{Tablas}
        \tiny
        \begin{itemize}\setlength{\itemsep}{-0.05cm}
          \item Título \azul{arriba}, numerado: ``Tabla 1. Distribución del
                gasto mensual''
          \item Unidades en el encabezado, no en cada celda
          \item Mismo número de decimales en toda la columna
          \item Fila de totales
          \item Fuente al pie
        \end{itemize}
      \end{block}
    \end{column}
    \begin{column}{0.48\textwidth}
      \begin{block}{Gráficos}
        \tiny
        \begin{itemize}\setlength{\itemsep}{-0.05cm}
          \item Título \azul{abajo}, numerado: ``Figura 1. Histograma del
                gasto''
          \item Ejes rotulados con unidades
          \item Eje vertical desde cero
          \item Indicar $n$
          \item Sin efectos 3D
        \end{itemize}
      \end{block}
    \end{column}
  \end{columns}

  \vspace{0.05cm}
  \begin{alertblock}{La regla de la referencia cruzada}
    \tiny
    Toda tabla y todo gráfico debe estar \destacar{citado en el texto}:
    ``como muestra la Tabla 1\dots''. Si un elemento no se menciona,
    \azul{sobra}.
  \end{alertblock}
\end{frame}

\begin{frame}{Ejemplo de párrafo de resultados}
  \begin{exampleblock}{Redacción profesional}
    \scriptsize
    ``La Tabla 1 muestra la distribución del gasto mensual de los 40 clientes
    encuestados en agosto de 2026. El gasto promedio fue de
    \azul{227\,000 pesos} con una desviación estándar de 56\,340
    (CV $= 24{,}8\%$), lo que indica una clientela razonablemente
    homogénea. La mitad de los clientes gastó menos de \azul{224\,730
    pesos} y el rango intercuartílico fue de 81\,270, de modo que el 50\%
    central se movió entre 186\,290 y 267\,560 pesos.

    \vspace{0.1cm}
    El histograma (Figura 1) presenta un único máximo en el intervalo
    $[192, 232)$ y una \azul{cola hacia la derecha}, confirmada por un
    coeficiente de asimetría de $+0{,}121$. No se detectaron valores
    atípicos según el criterio de $1{,}5 \times RIC$.''
  \end{exampleblock}

  \vspace{0.1cm}
  \begin{alertblock}{Qué tiene este párrafo}
    \tiny
    Cada cifra con su unidad; cada afirmación anclada a una tabla o figura; el
    vocabulario técnico aparece \destacar{después} de la explicación en
    lenguaje común, no antes.
  \end{alertblock}
\end{frame}

\begin{frame}{Taller final — informe completo}
  \begin{alertblock}{En grupos de 3 (35 minutos)}
    \scriptsize
    Con \destacar{todos} los resultados de la Cafetería El Portal calculados en
    las sesiones 5, 7, 8 y 9, redacte un informe estadístico completo.

    \vspace{0.12cm}
    \centering
    \begin{tabular}{@{}lrlr@{}}
      \toprule
      $n$ & 45 & $s$ & 8{,}28 \\
      $\bar{x}$ & 22{,}33 & $CV$ & 37{,}1\% \\
      $Me$ & 21{,}92 & $RIC$ & 11{,}88 \\
      $Mo$ & 21{,}43 & $As_2$ & $+0{,}149$ \\
      $Q_1$ / $Q_3$ & 16{,}17 / 28{,}05 & Zona principal & Centro (37,8\%) \\
      \bottomrule
    \end{tabular}

    \vspace{0.15cm}
    \raggedright
    Debe incluir: objetivo, metodología, tabla de variables, una tabla de
    resultados, la lectura del histograma, \destacar{tres} hallazgos
    interpretados, \destacar{dos} limitaciones y \destacar{dos}
    recomendaciones ancladas a cifras. Máximo dos páginas.
  \end{alertblock}
\end{frame}

\begin{frame}{Rúbrica: así se evalúa un informe}
  \scriptsize
  \begin{center}
  \begin{tabular}{@{}p{3.4cm}rp{6.6cm}@{}}
    \toprule
    \textbf{Criterio} & \textbf{\%} & \textbf{Qué se busca} \\
    \midrule
    Corrección de los cálculos & 25 &
      Las cifras coinciden con los datos; fórmulas bien aplicadas \\
    Elección de medidas y gráficos & 20 &
      Coherentes con el tipo y la escala de la variable \\
    Interpretación & 25 &
      Frases completas, con unidad y contexto; sin sobreinterpretar \\
    Limitaciones & 10 &
      Reconoce el alcance real de la muestra \\
    Recomendaciones & 10 &
      Concretas y ancladas a una cifra del estudio \\
    Presentación & 10 &
      Tablas y figuras numeradas, citadas y con fuente \\
    \bottomrule
  \end{tabular}
  \end{center}

  \vspace{0.05cm}
  {\tiny\color{uniGris} Esta misma rúbrica orienta la parte aplicada del
  Examen 3 (Unidad 3).}
\end{frame}

\begin{frame}[fragile]{Comparar e interpretar, en R}
  \scriptsize
\begin{lstlisting}[style=Rmini]
# Una medida por grupo: la comparacion de la lamina anterior, en una linea
tapply(gasto, sucursal, mean)
tapply(gasto, sucursal, median)
round(tapply(gasto, sucursal, sd) / tapply(gasto, sucursal, mean) * 100, 1)

# Todo el resumen por grupo
by(gasto, sucursal, summary)

# Tendencia: variacion porcentual mes a mes
ventas <- c(120, 128, 135, 132, 145, 158)
round(diff(ventas) / head(ventas, -1) * 100, 1)
round((tail(ventas, 1) / ventas[1] - 1) * 100, 1)   # variacion total
\end{lstlisting}
\begin{lstlisting}[style=Rout]
[1]  6.7  5.5 -2.2  9.8  9.0
[1] 31.7
\end{lstlisting}

  \vspace{-0.1cm}
  \begin{alertblock}{Lo que R no hace}
    \tiny
    R entrega las cifras en segundos, pero \destacar{no escribe la
    interpretación}, no reconoce una muestra sesgada y no distingue asociación
    de causa. Esos tres juicios siguen siendo suyos, y son justamente lo que se
    evalúa en esta unidad.
  \end{alertblock}
\end{frame}

\begin{frame}{Cierre de la sesión}
  \begin{block}{Lo que debe quedar claro hoy}
    \scriptsize
    \begin{itemize}\setlength{\itemsep}{0.02cm}
      \item Una interpretación necesita \azul{qué se midió, la cifra con
            unidad, su significado y una comparación}.
      \item Se buscan cinco patrones: concentración, dispersión, asimetría,
            atípicos y bimodalidad.
      \item Las tendencias se leen con \azul{variación porcentual} y con
            comparaciones entre grupos.
      \item Asociación no es causa, y el promedio puede no describir a nadie.
      \item El informe tiene una estructura fija; toda tabla y figura va
            numerada, citada y con fuente.
    \end{itemize}
  \end{block}

  \vspace{0.1cm}
  \begin{alertblock}{Trabajo independiente}
    \scriptsize
    Terminar el informe de El Portal y resolver los 20 ejercicios siguientes.
    La sesión 11 es el \destacar{repaso integrador de la Unidad 2} y el
    simulacro del Examen 2.
  \end{alertblock}
\end{frame}

\begin{frame}{Ejercicios propuestos}
  \vspace{-0.1cm}
  \begin{beamercolorbox}[wd=\textwidth,sep=0.10cm]{block body}
    \centering\scriptsize
    \textbf{1--7:} interpretación de cifras. \quad \textbf{8--14:} tendencias
    y comparación. \quad \textbf{15--20:} informe.
  \end{beamercolorbox}
  \listatiny
  \begin{columns}[T,onlytextwidth]
    \begin{column}{0.47\textwidth}
      \begin{enumerate}\setlength{\itemsep}{0.16cm}
        \item Interprete $\bar{x} = 3{,}4$ hijos por hogar
        \item Interprete $Me = 2$ en ese mismo estudio
        \item ¿Qué asimetría sugieren esas dos cifras juntas?
        \item Interprete $CV = 12\%$
        \item Interprete $CV = 68\%$
        \item Interprete $Q_3 = 45$ minutos
        \item Interprete $g_2 = +1{,}8$
        \item Ventas: 200 y 230. Variación porcentual
        \item Ventas: 230 y 207. Variación porcentual
        \item De 200 a 260 en dos periodos: variación total
      \end{enumerate}
    \end{column}
    \begin{column}{0.50\textwidth}
      \begin{enumerate}\setlength{\itemsep}{0.16cm}
        \setcounter{enumi}{10}
        \item Dos grupos, misma media y distinto $CV$: ¿qué concluye?
        \item Dos grupos, misma mediana y distinto máximo: ¿qué revisa?
        \item ¿Por qué comparar medianas y no solo medias?
        \item Un gráfico sube fuerte pero el eje empieza en 95: ¿qué hace?
        \item ¿Qué va en la sección de metodología?
        \item ¿Qué va en limitaciones?
        \item ¿Dónde se ubica el título de una tabla? ¿Y el de una figura?
        \item ¿Por qué el resumen se escribe de último?
        \item Escriba una recomendación anclada a $Q_3 = 267{,}56$
        \item ¿Qué diferencia hay entre un hallazgo y una recomendación?
      \end{enumerate}
    \end{column}
  \end{columns}
\end{frame}

\begin{frame}{Respuestas}{Autoevaluación de los ejercicios propuestos}
  \listatiny
  \begin{columns}[T,onlytextwidth]
    \begin{column}{0.47\textwidth}
      \begin{enumerate}\setlength{\itemsep}{0.16cm}
        \item En promedio cada hogar tiene 3,4 hijos
        \item La mitad de los hogares tiene 2 hijos o menos
        \item $\bar{x} > Me$: asimetría positiva
        \item Grupo muy homogéneo
        \item Grupo muy disperso; la media describe mal
        \item El 75\% espera menos de 45 minutos
        \item Leptocúrtica: pico alto y colas pesadas
        \item $+15{,}0\%$
        \item $-10{,}0\%$
        \item $+30{,}0\%$
      \end{enumerate}
    \end{column}
    \begin{column}{0.50\textwidth}
      \begin{enumerate}\setlength{\itemsep}{0.16cm}
        \setcounter{enumi}{10}
        \item Mismo centro, distinta consistencia
        \item Posibles atípicos o cola larga en uno de ellos
        \item Porque la mediana no se deja arrastrar por los extremos
        \item Rehacerlo con el eje desde cero
        \item Población, muestra, $n$, periodo e instrumento
        \item El alcance real: tamaño, periodo y posibles sesgos
        \item Tabla: arriba. Figura: abajo
        \item Porque solo se resume lo que ya está escrito
        \item Respuesta abierta: p. ej. programa de fidelización para el
              cuartil superior
        \item El hallazgo describe; la recomendación propone una acción
      \end{enumerate}
    \end{column}
  \end{columns}
\end{frame}

\begin{frame}{Referencias}
  \footnotesize
  \begin{itemize}
    \item[{[1]}] Monroy Saldívar, S. (2008). \textit{Estadística descriptiva}.
      Instituto Politécnico Nacional.
      \texttt{https://elibro.net/es/lc/uniminuto/titulos/74722}
    \item[{[2]}] Ordóñez Fernández, F. F. \& González Fernández, J. (2021).
      \textit{Estadística descriptiva paso a paso}. Ediciones Pirámide.
      \texttt{https://elibro.net/es/lc/uniminuto/titulos/215449}
    \item[{[3]}] Llinás Solano, H. (2014). \textit{Introducción a la
      estadística con aplicaciones en ciencias sociales}. Universidad del
      Norte. \texttt{https://elibro.net/es/lc/uniminuto/titulos/70062}
  \end{itemize}

  \vspace{0.35cm}
  \begin{center}
    {\Large\color{uniAzul}\textbf{¿Preguntas?}}
  \end{center}
\end{frame}

\end{document}
